%%%%%%%%%%%%%%%%%%%%%%%%%%%%%%%%%%%%%%%%%%%%%%%%%%%%%%%%%%%%%%%%%%%%
%% I, the copyright holder of this work, release this work into the
%% public domain. This applies worldwide. In some countries this may
%% not be legally possible; if so: I grant anyone the right to use
%% this work for any purpose, without any conditions, unless such
%% conditions are required by law.
%%%%%%%%%%%%%%%%%%%%%%%%%%%%%%%%%%%%%%%%%%%%%%%%%%%%%%%%%%%%%%%%%%%%

\documentclass{beamer}
\usetheme[faculty=fi]{fibeamer}
\usepackage[T2A]{fontenc}
\usepackage[utf8]{inputenc}
\usepackage[
  main=ukrainian, %% By using `czech` or `slovak` as the main locale
                %% instead of `english`, you can typeset the
                %% presentation in either Czech or Slovak,
                %% respectively.
  english, russian %% The additional keys allow foreign texts to be
]{babel}        %% typeset as follows:
%%
%%   \begin{otherlanguage}{czech}   ... \end{otherlanguage}
%%   \begin{otherlanguage}{slovak}  ... \end{otherlanguage}
%%
%% These macros specify information about the presentation
\title{Розробка моделей відкритого тексту на основі стохастичних контекстно-вільних граматик в нормальній формі Грейбах} %% that will be typeset on the
%\subtitle{Presentation Subtitle}
\subtitle{Грубіян Євген Олександрович}

\author{Керівник: к.ф.-м.н. Фесенко Андрій В'ячеславович}

%% These additional packages are used within the document:
\usepackage{ragged2e}  % `\justifying` text
\usepackage{booktabs}  % Tables
\usepackage{tabularx}
\usepackage{tikz}      % Diagrams
\usetikzlibrary{calc, shapes, backgrounds}
\usepackage{amsmath, amssymb}
\usepackage{url}       % `\url`s
\usepackage{listings}  % Code listings
\frenchspacing
\begin{document}
  \frame{\maketitle}


  \begin{darkframes}
    \section{Dark Frames}
    \subsection{Вступ}
    \begin{frame}{Актуальність}
      Мови та граматики моделюють складні внутрішні структури багатьох об'єктів із різних галузей науки.
      \begin{itemize}
        \item Формалізація природної мови - одна із найважливіших задач штучного інтелекту[Jurafsky, 2009].
        \item В криптоаналізі вдала модель відкритого тексту дає криптоаналітику додаткову інформацію для зламу[Shannon, 1948; Яглом, 1973].
        \item Структури послідовностей нуклеотидів в ДНК та РНК утворюються за граматичними правилами[Sakakibara, 1994]
      \end{itemize}
    \end{frame}

    \begin{frame}{...Актуальність}
      \framesubtitle{Jabberwocky, Льюїс Керрол}%
      \begin{tikzpicture}[overlay,remember picture]
        \node[anchor=south east,xshift=-30pt,yshift=35pt]
          at (current page.south east) {
            \includegraphics[width=35mm]{resources/jabberwocky-dark}
          };
      \end{tikzpicture}%
      'Twas brillig, and the slithy toves\\
      Did gyre and gimble in the wabe;\\
      All mimsy were the borogoves,\\
      And the mome raths outgrabe.\\\bigskip

      “Beware the Jabberwock, my son!\\
      The jaws that bite, the claws that catch!\\
      Beware the Jubjub bird, and shun\\
      The frumious Bandersnatch!”\\
    \end{frame}

    \begin{frame}[label=lists]{Мета, об'єкт та предмет дослідження}
      \begin{exampleblock}{Мета}
      Побудова нової моделi природних мов, що дасть змогу краще використовувати структуру мов для методiв криптоаналiзу
      \end{exampleblock}
      \begin{block}{Об'єкт}
        Формальна модель природних мов
      \end{block}
      \begin{block}{Предмет}
        Побудова полiпшеної моделi природних мов на основi стохастичної контекстно-вiльної граматики в нормальнiй формi Грейбах
      \end{block}

    \end{frame}

    \begin{frame}{Завдання}
      \begin{enumerate}
      \item Провести огляд опублікованих робіт за тематикою дослідження
      \item Розробити модель стохастичних регулярних граматик на основі прихованих моделей Маркова
      \item Побудувати модель стохастичних контекстно-вільних(КВ) граматик на основі узагальнення прихованих моделей Маркова та дослідити властивості цієї моделі
      \item Розробити програмну реалізацію моделей
      \item Перевірити узгодженість моделей емпіричним шляхом
      \end{enumerate}
    \end{frame}

    \begin{frame}{Моделі відкритого тексту}
      \begin{itemize}
        \item Класичні моделі M0, M1, M2, M3
        \item n-грамна модель
        \item Позиційні моделі
        \item Моделі на основі нейронних мереж
        \item Граматичні моделі
      \end{itemize}
      \begin{exampleblock}{Визначення 1.7}
      \textit{Контекстно-вільною граматикою} $G$ називається граматика $ G = \big\langle N, \Sigma, S, R \big\rangle $, всі правила виводу $R$ якої мають вигляд:
      $$ A \rightarrow \gamma,\ A \in N,\ \gamma \in \left(N \cup \Sigma \right)^* $$
      \end{exampleblock}
    \end{frame}

    \subsection{Результати}
    \begin{frame}{Модель стохастичної регулярної граматики}
      \begin{block}{Теорема 2.1}
        Для кожної стохастичної регулярної граматики iснує прихована модель Маркова другого порядку за спостереженнями, що допускає таку саму стохастичну мову.
      \end{block}
      \begin{block}{Твердження 2.1}
        Алгоритм перетворення стохастичної регулярної граматики до прихованої моделi Маркова другого порядку за спостереженнями має складнiсть $\mathcal{O}(n)$, де $n$ - кiлькiсть правил виводу в
        граматицi.
      \end{block}
    \end{frame}

    \begin{frame}{Модель стохастичної КВ граматики}
      \begin{exampleblock}{Визначення 2.4}
        \textit{Прихованою моделлю Маркова другого порядку за спостереженнями зі стеком} з простором латентних(прихованих) станів $E = \{ 1,\dots,n \}$, простором спостережень $O = \{ 1,\dots,m \}$ та стеком  називається стохастичний процес $ \{ (X_t, S_t, Y_t),\ t \in \mathbb{N} \},\ \ X_t \in E,\ Y_t \in O $, $S_t$ - стек над множиною $E$, переходи в якому здійснюються за алгоритмом:
        \begin{itemize}
          \item Ініціалізація стеку: $S_1 = \emptyset,\ S_t = S_{t-1}$

        \end{itemize}
      \end{exampleblock}
    \end{frame}

    \begin{frame}{Модель стохастичної КВ граматики}
      \begin{itemize}
        \item Перехід між латентними станами:
        \begin{multline}
          P(X_t = i,\ S_t.Push(k) | X_{1:t-1}) = \\
          P(X_t = i,\ S_t.Push(k) | X_{t-1}),\\
          P(X_t = i,\ S_t.Push(k) | X_{t-1} = j) = p_{jik}
        \end{multline}
        \item Спостереження:
        \begin{multline}
          P(Y_t = i | X_{1:t}, S_t.Top ) = P(Y_t = i | X_t, X_{t-1}, S_t.Top),\\
          P(Y_t = i | X_t = j, X_{t-1} = k, S_t.Top = l) = q_{lkji}
        \end{multline}
        \item Якщо $X_t = \epsilon$ та $ S_t.Top \neq \epsilon $, то $ X_t := S_t.Top,\ S_t.Pop $. Якщо $X_t = \epsilon$ та $ S_t.Top = \epsilon $, то ланцюг завершується.
      \end{itemize}
    \end{frame}

    \begin{frame}{Модель стохастичної КВ граматики}
      \begin{block}{Теорема 2.2}
        Для кожної стохастичної контекстно-вільної граматики(\textit{SCFG}) існує прихована модель Маркова другого порядку за спостереженнями зі стеком(\alert{HMM2S}), що допускає таку саму стохастичну мову.
      \end{block}
      \begin{block}{Твердження 2.2}
      Алгоритм перетворення стохастичної граматики в 2-нормальній формі Грейбах до прихованої моделі Маркова другого порядку за спостереженнями зі стеком має складність $\mathcal{O}(n)$, де $n$ - кількість правил виводу в граматиці.
      \end{block}

    \end{frame}

    \begin{frame}{Узгодженість моделей}
      Було перевірено емпіричним шляхом узгодженість запропонованої моделі і стохастичних КВ граматик за допомогою розробленого програмного пакету: \url{https://github.com/juja256/hmm2s}
      \begin{alertblock}{Тестова граматика}
        \begin{columns}[onlytextwidth]
          \column{.5\textwidth}
          \begin{align*}
            & (0.2) A_1 \rightarrow a A_3 A_2 \\
            & (0.8) A_1 \rightarrow c A_2 \\
            & (0.5) A_2 \rightarrow a \\
            & (0.5) A_2 \rightarrow b
          \end{align*}
          \column{.5\textwidth}
          \begin{align*}
            & (0.3) A_3 \rightarrow d A_2 \\
            & (0.3) A_3 \rightarrow c A_2 \\
            & (0.4) A_3 \rightarrow a
          \end{align*}
        \end{columns}

      \end{alertblock}
    \end{frame}


    \begin{frame}{Узгодженість моделей}
      %\framesubtitle{Частотні таблиці}%
      \begin{table}[!b]
        {\carlitoTLF % Use monospaced lining figures
        \begin{tabularx}{\textwidth}{Xcc}
          \textbf{Слово} & \textbf{Частота / SCFG} & \textbf{Частота / HMM2S} \\
          \toprule
          a a a & 5 & 1 \\
          a a b & 1 & 2 \\
          a c a a & 1 & 1 \\
          a c a b & 3 & 0 \\
          a c b a & 2 & 0 \\
          a c b b & 3 & 2 \\
          a d a a & 0 & 2 \\
          a d a b & 0 & 2 \\
          a d b a & 1 & 1 \\
          a d b b & 1 & 3 \\
          c a & 45 & 45 \\
          c b & 38 & 41
          %\bottomrule
        \end{tabularx}}

      \end{table}

    \end{frame}

  \end{darkframes}


  \section{Light Frames}
    \subsection{Висновки}
    \begin{frame}{Висновки}
      \begin{itemize}
        \item Розроблено нову модель природної мови(відкритого тексту) на основі запропонованої прихованої моделі Маркова другого порядку за спостереженнями зі стеком
        \item Доведено еквівалентність запропонованої моделі та стохастичних КВ граматик та запропоновано поліноміальний алгоритм зведення стохастичної КВ граматики до запропонованої моделі
        \item Емпірично підтверджено узгодженість моделі і стохастичної КВ граматики
        \item Нова модель дозволяє розробку ефективніших алгоритмів навчання стохастичних КВ граматик
      \end{itemize}
    \end{frame}

\end{document}
